% Source-aligned English for printed p.209.
\clearpage
\SourcePage{209}

\LayerTarget{tohoku:conclusion:theorem-5.5.3:p1}
\noindent\textit{is injective if $n=1$.} Moreover:

\LayerTarget{tohoku:corollary:1-theorem-5.5.3}
\textsc{Corollary 1.} \textit{If all the $U_{i_0\ldots i_p}$ are
$A$-acyclic, then the homomorphisms \textup{(5.5.5)} are isomorphisms.}

\LayerTarget{tohoku:corollary:2-theorem-5.5.3}
\textsc{Corollary 2.} \textit{The preceding theorem remains valid if
${\mathbf U}$ is a closed $G$-covering, provided that one of the following two
conditions holds: a) ${\mathbf U}$ is locally finite and $X$ is paracompact;
b) ${\mathbf U}$ is finite.}

\LayerTarget{tohoku:proof:corollary-2-theorem-5.5.3}
In case a), by what is supposed to have been seen in 3.8 (see the remark in
3.8), conditions (i) and (ii) of Lemma 5.5.1 are still satisfied (for (i), this
is in fact trivial without any paracompactness hypothesis, since ${\mathbf U}$
is locally finite), and the formula
$R^q\mathbf K(A)=C({\mathbf U},H^q(A))$ remains valid. In case b), one must
follow Godement's method (which he develops for $G$ reduced to the identity),
which we shall sketch because of its possible applications to abstract
algebraic geometry. First, following an idea of P. Cartier, introduce the
complex of sheaves ${\mathbf C}({\mathbf U},A)$, whose complex of sections over any open
set $V$ is by definition $C({\mathbf U}|V,A|V)$ (the sign $|$ indicating, as
usual, restriction to the set written after it). This is proved to be a
\emph{resolution} of $A$ (see Godement's book [9] for details). We apply to it the
corollary to Lemma 5.5.2 (the letters in that lemma having the same meaning as
above). It remains only to prove that
$R^qF(\mathrm C)=H^q(X,{\mathbf C}({\mathbf U},A))$ is functorially isomorphic to
$C({\mathbf U},H^q(A))$. This presents no difficulty because ${\mathbf U}$ is
\emph{finite}, using Theorem 3.5.1.

\LayerTarget{tohoku:definition:fixed-point-free-g-cover}
We shall say that a \emph{$G$-covering} ${\mathbf U}=(U_i)_{i\in I}$ is
``\emph{fixed-point-free}'' if $G$ acts on $I$ without fixed points, i.e. if $g\ne e$
implies $g.\, i\ne i$ for every $i$. If $P$ is any $G$-presheaf, then for
every $n$ one has

\LayerTarget{tohoku:equation:5.5.6}
\[
\tag{5.5.6}
C^n({\mathbf U},P)=
\prod_{(i_0,\ldots,i_n)}\ \prod_{g\in G}
P(g.\, U_{i_0\ldots i_n}),
\]

\LayerTarget{tohoku:calculation:fixed-point-free-g-cover-cochains}
where the first sign $\displaystyle\prod$ extends over a complete system of representatives
$(i_0,\ldots,i_n)$ of $I^{n+1}/G$. In the second product all terms are
canonically isomorphic to $P(U_{i_0\ldots,i_n})$, so this product identifies
with the group of maps from $G$ to the fixed group $P(U_{i_0\ldots i_n})$, on
which $G$ acts by left translations. Such a $G$-module is well known to be
$\Gamma^G$-acyclic, and hence so is $C^n({\mathbf U},P)$. One must note,
however, that formula (5.5.6) is valid only when $C({\mathbf U},P)$ denotes the
complex of \emph{arbitrary} cochains (not necessarily alternating) of ${\mathbf U}$
with coefficients in $P$ (whereas the preceding considerations applied
equally to arbitrary or alternating cochains). Using the first hyperhomology
spectral sequence for the functor $\Gamma^G$ relative to the complex
$C({\mathbf U},P)$, one obtains:

\LayerTarget{tohoku:proposition:5.5.4}
\textsc{Proposition 5.5.4.} \textit{If ${\mathbf U}$ is a
fixed-point-free $G$-covering, then for every $G$-presheaf $P$ one has:}

\LayerTarget{tohoku:equation:5.5.7}
\[
\tag{5.5.7}
H^n({\mathbf U},G,P)=H^n\bigl((C({\mathbf U},P))^G\bigr)
\]

\TohokuContinuation{210}
